\section{Forecasting features from local reconstruction}
\label{sec:reconstruction-method}
\label{sec:reconstruction-mechanism}

Let $(y_t)$ be a scalar observed series. At forecast origin $t$, the
available history is $x_t=(y_{t-d+1},\ldots,y_t)^\top\in\R^d$,
and the target at horizon $h\geq1$ is $y_{t+h}$. A forecasting procedure
maps the history to a separate prediction for each $h$ in a specified
finite horizon set $\mathcal H$. The representation and decoder are
fitted using a set A of history--target pairs. Distinct episode sets B,
C, and D are used for model selection, correction calibration, and
final evaluation, respectively; Section~\ref{sec:experiment} specifies
their construction.

Vectors are columns. The notation $v^\top$ denotes transpose,
$\|v\|_2$ the Euclidean norm, and $\|v\|_1$ the sum of absolute
coordinates. The matrix $I_p$ is the $p\times p$ identity, and
$\mathbf1_n$ is the length-$n$ vector of ones; omitted subscripts are
determined by the displayed products. Expectations and variances refer
to the probability law specified for each statement or experiment.
All feature-space statements condition on a fixed fitted reference bank.

\subsection{Reference banks and reconstruction maps}


Let $z\in\R^d$ be the coordinatewise standardized version of a history
$x$, using means and scales estimated on A. A local reference bank is a collection of fitted models indexed by
$j$. Model $j$ has $n_j$ reference histories stored as the rows of
$Q_j\in\R^{n_j\times d}$, center $\mu_j\in\R^d$, and orthonormal
principal directions $P_j\in\R^{d\times r_j}$. Here $r_j$ is its local
rank and $P_j^\top P_j=I_{r_j}$. Write $q_{ji}\in\R^d$ for reference
$i$, so row $i$ of $Q_j$ is $q_{ji}^\top$. Let $S_j$ be the sample
covariance matrix of the projected references $P_j^\top q_{ji}$. For
positive rank and positive-definite $S_j$, the shared Gaussian kernel is
\[
 \kappa_j(z,z')=\exp\!\left\{
 -\frac{[P_j^\top(z-z')]^\top S_j^{-1}[P_j^\top(z-z')]}{2r_j}
 \right\}.
\]
Define the kernel vector $k_j(z)\in\R^{n_j}$ by
$[k_j(z)]_i=\kappa_j(z,q_{ji})$, and the Gram matrix
$K_j\in\R^{n_j\times n_j}$ by
$[K_j]_{i\ell}=\kappa_j(q_{ji},q_{j\ell})$.
With a positive reconstruction regularizer $\lambda_j$, set
$C_j=(K_j+\lambda_j I_{n_j})^{-1}$. The Gaussian Gram matrix and its
regularized inverse $C_j$ are symmetric. The regularizer is fitted from A
by the reference-reconstruction rule specified in
Appendix~\ref{app:reconstruction-results}.

The KAHM and matched positive-kernel reconstructions are
\begin{equation}
 A_{K,j}(z)=\frac{Q_j^\top C_j k_j(z)}
 {\mathbf1^\top C_j k_j(z)},\qquad
 A_{R,j}(z)=\frac{Q_j^\top k_j(z)}{\mathbf1^\top k_j(z)}.
 \label{eq:mechanism-main-maps}
\end{equation}
Both take values in $\R^d$. The first is defined when its signed
denominator is nonzero; the second has a positive denominator for finite
inputs. The matched affine reconstruction is
$A_{A,j}(z)=\mu_j+P_jP_j^\top(z-\mu_j)$.
\subsection{Common features and reconstruction departures}

Models are grouped into partitions. If $\mathcal J_b$ indexes the models
in partition $b$, their shared locality weights are
\[
 g_j(z)=\frac{\exp\{-\|z-\mu_j\|_2^2/(2s_j)\}}
 {\sum_{\ell\in\mathcal J_b}
       \exp\{-\|z-\mu_\ell\|_2^2/(2s_\ell)\}},
 \qquad j\in\mathcal J_b,
\]
where $s_j>0$ is the fitted local squared-radius scale. These weights,
called gates, sum to one within each partition and do not depend on the
reconstruction being compared.

The common feature vector $\psi_0(z)$ contains all gates $g_j(z)$ and
all weighted principal-coordinate vectors $g_j(z)P_j^\top(z-\mu_j)$.
Define the reconstruction departures by
\begin{equation}
 \begin{split}
 d_K(z)&=\operatorname{concat}_{j}
       \bigl[g_j(z)\{A_{K,j}(z)-A_{A,j}(z)\}\bigr],\\
 d_R(z)&=\operatorname{concat}_{j}
       \bigl[g_j(z)\{A_{R,j}(z)-A_{A,j}(z)\}\bigr].
 \end{split}
 \label{eq:mechanism-main-departures}
\end{equation}
The operator $\operatorname{concat}$ stacks its column-vector arguments
in a fixed order. Each departure retains all history coordinates,
including components outside the local principal subspace. Substituting
the affine map for the nonlinear reconstruction makes the departure
identically zero. The common-feature procedure is therefore the nested
control without reconstruction departures. The RBF departure is matched
to the KAHM departure in references, kernels, gates, and feature dimension.


The affine control is nested at the level of the raw feature construction.
The comparison with positive-kernel RBF reconstruction additionally holds
fixed the number and locations of references, the projected kernel, the
locality weights, and the decoder architecture. The difference between
the two raw departure blocks is therefore the use of $C_j$ in the
normalized reconstruction. This matching does not assert equality of
fitted predictions, and it does not compare KAHM with every possible
kernel or nonlinear decoder.

\subsection{Defined evaluation of signed normalizations}

To guard against cancellation in the signed KAHM normalization, define
\begin{equation}
 a_j(z)=C_j k_j(z),\qquad
 \rho_j(z)=\frac{|\mathbf1^\top a_j(z)|}{\|a_j(z)\|_1}
 \quad\text{when }\|a_j(z)\|_1>0.
 \label{eq:mechanism-relative-denominator}
\end{equation}
The ratio is invariant to a common positive rescaling of the query kernel.
When $\rho_j(z)\le10^{-10}$, or a required quantity is undefined or
nonfinite, the local reconstruction is replaced by its affine map, making
that raw departure zero. The same finite-value checks apply to the positive
RBF map. All affected rows are recorded. This fallback gives a defined
feature vector where the original ratio cannot be evaluated reliably;
it does not establish continuity across the fallback boundary.


A zero raw departure under fallback need not remain zero after A-fitted
feature centering. The fallback defines the input to the subsequent
transformation; the norm bound in Section~\ref{sec:bounded-decoder}
applies to that transformed input. The algebraic propositions in
Section~\ref{sec:reconstruction-theory} concern the exact ratios on their stated
domains, rather than the piecewise numerical fallback procedure.
